% 决赛策略研究报告 公共导言（XeLaTeX + ctexart）
% 原则：层级清楚（章、节两级为主，最多到小节）；正文成段；不用彩色框、侧栏、底纹、小字注释。
% 图表标题在上，写成「图 1　标题（数据范围）」，数据来源写进标题括号或正文，不另起小字行。
%
% 用法：
%   \documentclass[UTF8,a4paper,zihao=-4,fontset=none]{ctexart}
%   % 决赛策略研究报告 公共导言（XeLaTeX + ctexart）
% 原则：层级清楚（章、节两级为主，最多到小节）；正文成段；不用彩色框、侧栏、底纹、小字注释。
% 图表标题在上，写成「图 1　标题（数据范围）」，数据来源写进标题括号或正文，不另起小字行。
%
% 用法：
%   \documentclass[UTF8,a4paper,zihao=-4,fontset=none]{ctexart}
%   % 决赛策略研究报告 公共导言（XeLaTeX + ctexart）
% 原则：层级清楚（章、节两级为主，最多到小节）；正文成段；不用彩色框、侧栏、底纹、小字注释。
% 图表标题在上，写成「图 1　标题（数据范围）」，数据来源写进标题括号或正文，不另起小字行。
%
% 用法：
%   \documentclass[UTF8,a4paper,zihao=-4,fontset=none]{ctexart}
%   % 决赛策略研究报告 公共导言（XeLaTeX + ctexart）
% 原则：层级清楚（章、节两级为主，最多到小节）；正文成段；不用彩色框、侧栏、底纹、小字注释。
% 图表标题在上，写成「图 1　标题（数据范围）」，数据来源写进标题括号或正文，不另起小字行。
%
% 用法：
%   \documentclass[UTF8,a4paper,zihao=-4,fontset=none]{ctexart}
%   \input{../_style/report_preamble.tex}
%   \title{…} \subtitle{…} \team{…} \reportdate{…}
%   \begin{document} \makereporttitle ...

\setCJKmainfont{Songti SC}[BoldFont={Songti SC Bold}]
\def\yhpath{/Applications/Microsoft PowerPoint.app/Contents/Resources/DFonts/}
\setCJKsansfont{msyh.ttc}[Path=\yhpath,BoldFont=msyhbd.ttc]
\setCJKmonofont{msyh.ttc}[Path=\yhpath]
\setmainfont{Times New Roman}
\setsansfont{msyh.ttc}[Path=\yhpath,BoldFont=msyhbd.ttc]
\setmonofont{Menlo}[Scale=0.88]

\usepackage{newunicodechar}
\newunicodechar{→}{\ensuremath{\rightarrow}}
\newunicodechar{≈}{\ensuremath{\approx}}
\newunicodechar{≤}{\ensuremath{\leq}}
\newunicodechar{≥}{\ensuremath{\geq}}
\newunicodechar{−}{\ensuremath{-}}
\newunicodechar{×}{\ensuremath{\times}}
\newunicodechar{÷}{\ensuremath{\div}}
\newunicodechar{±}{\ensuremath{\pm}}
\newunicodechar{β}{\ensuremath{\beta}}
\newunicodechar{λ}{\ensuremath{\lambda}}
\newunicodechar{σ}{\ensuremath{\sigma}}

\usepackage{xcolor}
\definecolor{rnavy}{HTML}{1B365D}
\definecolor{rgrey}{HTML}{7A828C}
\definecolor{rline}{HTML}{BFC5CC}
\definecolor{raccent}{HTML}{C0392B}

\usepackage[a4paper,top=2.6cm,bottom=2.6cm,left=2.7cm,right=2.7cm,headheight=14pt]{geometry}
\linespread{1.45}
\setlength{\parskip}{0.35em}

% 标题层级：章（\section）、节（\subsection）、小节（\subsubsection，少用）
\ctexset{
  section={format=\sffamily\Large\bfseries\color{rnavy}\raggedright, name={,}, number=\arabic{section},
           aftername=\enspace, beforeskip=1.6em, afterskip=0.8em,
           numberformat={\color{white}\colorbox{rnavy}{\makebox[1.1em]{\arabic{section}}}}},
  subsection={format=\sffamily\large\bfseries\color{rnavy}\raggedright, beforeskip=1.2em, afterskip=0.5em},
  subsubsection={format=\sffamily\normalsize\bfseries\raggedright, beforeskip=0.9em, afterskip=0.3em},
}

% 页眉：左栏目、右日期，下方细色线；页脚：页码
\usepackage{fancyhdr}
\providecommand{\ReportHeadLeft}{策略研究报告}
\providecommand{\ReportHeadRight}{2026 年 11 月}
\pagestyle{fancy}
\fancyhf{}
\fancyhead[L]{\sffamily\small\color{rnavy}\ReportHeadLeft}
\fancyhead[R]{\sffamily\small\color{rnavy}\ReportHeadRight}
\renewcommand{\headrule}{{\color{rnavy}\hrule height 0.8pt}}
\fancyfoot[C]{\sffamily\small\thepage}
\fancypagestyle{plain}{\fancyhf{}\renewcommand{\headrulewidth}{0pt}\fancyfoot[C]{\sffamily\small\thepage}}

% 图表：标题在上，「图 1　标题」
\usepackage{graphicx,float}
\usepackage[font=small,labelfont={sf,bf,color=rnavy},textfont={sf,bf,color=rnavy},labelsep=colon,justification=raggedright,singlelinecheck=false,skip=5pt]{caption}
\DeclareCaptionLabelSeparator{zhcolon}{：}
\captionsetup{labelsep=zhcolon}
\captionsetup[figure]{name=图表,position=top}
\captionsetup[table]{name=图表,position=top}
\makeatletter\let\c@table\c@figure\makeatother
% 图表下方一行资料来源（唯一允许的小字）
\newcommand{\source}[1]{\par\vspace{2pt}{\raggedright\sffamily\footnotesize\color{rgrey}资料来源：#1\par}}

\usepackage{booktabs,tabularx,array,multirow}
\renewcommand{\arraystretch}{1.3}
\newcolumntype{L}{>{\raggedright\arraybackslash}X}
\newcolumntype{C}{>{\centering\arraybackslash}X}

\usepackage{amsmath,amssymb,bm}
\usepackage{enumitem}
\setlist{nosep,leftmargin=2em,itemsep=0.2em}

\usepackage{tikz}
\usetikzlibrary{arrows.meta,positioning,calc,fit,decorations.pathreplacing}
\tikzset{
  node/.style = {draw=rnavy, line width=0.6pt, rounded corners=1.5pt, align=flush center,
                 inner xsep=6pt, inner ysep=4pt, font=\sffamily\small},
  key/.style  = {node, draw=raccent, line width=0.8pt},
  soft/.style = {node, draw=rline},
  flow/.style = {-{Stealth[length=2mm]}, line width=0.6pt, draw=rnavy},
  note/.style = {font=\sffamily\small, text=rgrey, align=flush center},
}

% 首页研究要点：每条以一句加粗结论开头
\newenvironment{keypoints}{%
  \par\noindent{\sffamily\large\bfseries\color{rnavy}研究要点}\par\vspace{0.3em}
  {\color{rnavy}\hrule height 0.6pt}\vspace{0.6em}
  \begin{itemize}[leftmargin=1.2em,itemsep=0.6em,label={\color{rnavy}\rule[0.3ex]{0.45em}{0.45em}}]}
  {\end{itemize}\vspace{0.4em}{\color{rline}\hrule height 0.4pt}}
\newcommand{\kp}[1]{{\sffamily\bfseries\color{rnavy}#1}}

\usepackage[hidelinks]{hyperref}
\usepackage{bookmark}

% 封面
\makeatletter
\newcommand{\subtitle}[1]{\gdef\@subtitle{#1}}
\newcommand{\team}[1]{\gdef\@team{#1}}
\newcommand{\reportdate}[1]{\gdef\@reportdate{#1}}
\providecommand{\@subtitle}{}\providecommand{\@team}{}\providecommand{\@reportdate}{}
\newcommand{\makereporttitle}{%
  \thispagestyle{fancy}
  \vspace*{0.4em}
  {\raggedright\sffamily\color{rnavy}\fontsize{20pt}{27pt}\selectfont\bfseries\@title\par}
  \vspace{0.5em}
  {\raggedright\sffamily\large\color{rnavy}\@subtitle\par}
  \vspace{0.3em}
  {\raggedright\sffamily\small\color{rgrey}\@team\qquad\@reportdate\par}
  \vspace{1.2em}}
\makeatother

%   \title{…} \subtitle{…} \team{…} \reportdate{…}
%   \begin{document} \makereporttitle ...

\setCJKmainfont{Songti SC}[BoldFont={Songti SC Bold}]
\def\yhpath{/Applications/Microsoft PowerPoint.app/Contents/Resources/DFonts/}
\setCJKsansfont{msyh.ttc}[Path=\yhpath,BoldFont=msyhbd.ttc]
\setCJKmonofont{msyh.ttc}[Path=\yhpath]
\setmainfont{Times New Roman}
\setsansfont{msyh.ttc}[Path=\yhpath,BoldFont=msyhbd.ttc]
\setmonofont{Menlo}[Scale=0.88]

\usepackage{newunicodechar}
\newunicodechar{→}{\ensuremath{\rightarrow}}
\newunicodechar{≈}{\ensuremath{\approx}}
\newunicodechar{≤}{\ensuremath{\leq}}
\newunicodechar{≥}{\ensuremath{\geq}}
\newunicodechar{−}{\ensuremath{-}}
\newunicodechar{×}{\ensuremath{\times}}
\newunicodechar{÷}{\ensuremath{\div}}
\newunicodechar{±}{\ensuremath{\pm}}
\newunicodechar{β}{\ensuremath{\beta}}
\newunicodechar{λ}{\ensuremath{\lambda}}
\newunicodechar{σ}{\ensuremath{\sigma}}

\usepackage{xcolor}
\definecolor{rnavy}{HTML}{1B365D}
\definecolor{rgrey}{HTML}{7A828C}
\definecolor{rline}{HTML}{BFC5CC}
\definecolor{raccent}{HTML}{C0392B}

\usepackage[a4paper,top=2.6cm,bottom=2.6cm,left=2.7cm,right=2.7cm,headheight=14pt]{geometry}
\linespread{1.45}
\setlength{\parskip}{0.35em}

% 标题层级：章（\section）、节（\subsection）、小节（\subsubsection，少用）
\ctexset{
  section={format=\sffamily\Large\bfseries\color{rnavy}\raggedright, name={,}, number=\arabic{section},
           aftername=\enspace, beforeskip=1.6em, afterskip=0.8em,
           numberformat={\color{white}\colorbox{rnavy}{\makebox[1.1em]{\arabic{section}}}}},
  subsection={format=\sffamily\large\bfseries\color{rnavy}\raggedright, beforeskip=1.2em, afterskip=0.5em},
  subsubsection={format=\sffamily\normalsize\bfseries\raggedright, beforeskip=0.9em, afterskip=0.3em},
}

% 页眉：左栏目、右日期，下方细色线；页脚：页码
\usepackage{fancyhdr}
\providecommand{\ReportHeadLeft}{策略研究报告}
\providecommand{\ReportHeadRight}{2026 年 11 月}
\pagestyle{fancy}
\fancyhf{}
\fancyhead[L]{\sffamily\small\color{rnavy}\ReportHeadLeft}
\fancyhead[R]{\sffamily\small\color{rnavy}\ReportHeadRight}
\renewcommand{\headrule}{{\color{rnavy}\hrule height 0.8pt}}
\fancyfoot[C]{\sffamily\small\thepage}
\fancypagestyle{plain}{\fancyhf{}\renewcommand{\headrulewidth}{0pt}\fancyfoot[C]{\sffamily\small\thepage}}

% 图表：标题在上，「图 1　标题」
\usepackage{graphicx,float}
\usepackage[font=small,labelfont={sf,bf,color=rnavy},textfont={sf,bf,color=rnavy},labelsep=colon,justification=raggedright,singlelinecheck=false,skip=5pt]{caption}
\DeclareCaptionLabelSeparator{zhcolon}{：}
\captionsetup{labelsep=zhcolon}
\captionsetup[figure]{name=图表,position=top}
\captionsetup[table]{name=图表,position=top}
\makeatletter\let\c@table\c@figure\makeatother
% 图表下方一行资料来源（唯一允许的小字）
\newcommand{\source}[1]{\par\vspace{2pt}{\raggedright\sffamily\footnotesize\color{rgrey}资料来源：#1\par}}

\usepackage{booktabs,tabularx,array,multirow}
\renewcommand{\arraystretch}{1.3}
\newcolumntype{L}{>{\raggedright\arraybackslash}X}
\newcolumntype{C}{>{\centering\arraybackslash}X}

\usepackage{amsmath,amssymb,bm}
\usepackage{enumitem}
\setlist{nosep,leftmargin=2em,itemsep=0.2em}

\usepackage{tikz}
\usetikzlibrary{arrows.meta,positioning,calc,fit,decorations.pathreplacing}
\tikzset{
  node/.style = {draw=rnavy, line width=0.6pt, rounded corners=1.5pt, align=flush center,
                 inner xsep=6pt, inner ysep=4pt, font=\sffamily\small},
  key/.style  = {node, draw=raccent, line width=0.8pt},
  soft/.style = {node, draw=rline},
  flow/.style = {-{Stealth[length=2mm]}, line width=0.6pt, draw=rnavy},
  note/.style = {font=\sffamily\small, text=rgrey, align=flush center},
}

% 首页研究要点：每条以一句加粗结论开头
\newenvironment{keypoints}{%
  \par\noindent{\sffamily\large\bfseries\color{rnavy}研究要点}\par\vspace{0.3em}
  {\color{rnavy}\hrule height 0.6pt}\vspace{0.6em}
  \begin{itemize}[leftmargin=1.2em,itemsep=0.6em,label={\color{rnavy}\rule[0.3ex]{0.45em}{0.45em}}]}
  {\end{itemize}\vspace{0.4em}{\color{rline}\hrule height 0.4pt}}
\newcommand{\kp}[1]{{\sffamily\bfseries\color{rnavy}#1}}

\usepackage[hidelinks]{hyperref}
\usepackage{bookmark}

% 封面
\makeatletter
\newcommand{\subtitle}[1]{\gdef\@subtitle{#1}}
\newcommand{\team}[1]{\gdef\@team{#1}}
\newcommand{\reportdate}[1]{\gdef\@reportdate{#1}}
\providecommand{\@subtitle}{}\providecommand{\@team}{}\providecommand{\@reportdate}{}
\newcommand{\makereporttitle}{%
  \thispagestyle{fancy}
  \vspace*{0.4em}
  {\raggedright\sffamily\color{rnavy}\fontsize{20pt}{27pt}\selectfont\bfseries\@title\par}
  \vspace{0.5em}
  {\raggedright\sffamily\large\color{rnavy}\@subtitle\par}
  \vspace{0.3em}
  {\raggedright\sffamily\small\color{rgrey}\@team\qquad\@reportdate\par}
  \vspace{1.2em}}
\makeatother

%   \title{…} \subtitle{…} \team{…} \reportdate{…}
%   \begin{document} \makereporttitle ...

\setCJKmainfont{Songti SC}[BoldFont={Songti SC Bold}]
\def\yhpath{/Applications/Microsoft PowerPoint.app/Contents/Resources/DFonts/}
\setCJKsansfont{msyh.ttc}[Path=\yhpath,BoldFont=msyhbd.ttc]
\setCJKmonofont{msyh.ttc}[Path=\yhpath]
\setmainfont{Times New Roman}
\setsansfont{msyh.ttc}[Path=\yhpath,BoldFont=msyhbd.ttc]
\setmonofont{Menlo}[Scale=0.88]

\usepackage{newunicodechar}
\newunicodechar{→}{\ensuremath{\rightarrow}}
\newunicodechar{≈}{\ensuremath{\approx}}
\newunicodechar{≤}{\ensuremath{\leq}}
\newunicodechar{≥}{\ensuremath{\geq}}
\newunicodechar{−}{\ensuremath{-}}
\newunicodechar{×}{\ensuremath{\times}}
\newunicodechar{÷}{\ensuremath{\div}}
\newunicodechar{±}{\ensuremath{\pm}}
\newunicodechar{β}{\ensuremath{\beta}}
\newunicodechar{λ}{\ensuremath{\lambda}}
\newunicodechar{σ}{\ensuremath{\sigma}}

\usepackage{xcolor}
\definecolor{rnavy}{HTML}{1B365D}
\definecolor{rgrey}{HTML}{7A828C}
\definecolor{rline}{HTML}{BFC5CC}
\definecolor{raccent}{HTML}{C0392B}

\usepackage[a4paper,top=2.6cm,bottom=2.6cm,left=2.7cm,right=2.7cm,headheight=14pt]{geometry}
\linespread{1.45}
\setlength{\parskip}{0.35em}

% 标题层级：章（\section）、节（\subsection）、小节（\subsubsection，少用）
\ctexset{
  section={format=\sffamily\Large\bfseries\color{rnavy}\raggedright, name={,}, number=\arabic{section},
           aftername=\enspace, beforeskip=1.6em, afterskip=0.8em,
           numberformat={\color{white}\colorbox{rnavy}{\makebox[1.1em]{\arabic{section}}}}},
  subsection={format=\sffamily\large\bfseries\color{rnavy}\raggedright, beforeskip=1.2em, afterskip=0.5em},
  subsubsection={format=\sffamily\normalsize\bfseries\raggedright, beforeskip=0.9em, afterskip=0.3em},
}

% 页眉：左栏目、右日期，下方细色线；页脚：页码
\usepackage{fancyhdr}
\providecommand{\ReportHeadLeft}{策略研究报告}
\providecommand{\ReportHeadRight}{2026 年 11 月}
\pagestyle{fancy}
\fancyhf{}
\fancyhead[L]{\sffamily\small\color{rnavy}\ReportHeadLeft}
\fancyhead[R]{\sffamily\small\color{rnavy}\ReportHeadRight}
\renewcommand{\headrule}{{\color{rnavy}\hrule height 0.8pt}}
\fancyfoot[C]{\sffamily\small\thepage}
\fancypagestyle{plain}{\fancyhf{}\renewcommand{\headrulewidth}{0pt}\fancyfoot[C]{\sffamily\small\thepage}}

% 图表：标题在上，「图 1　标题」
\usepackage{graphicx,float}
\usepackage[font=small,labelfont={sf,bf,color=rnavy},textfont={sf,bf,color=rnavy},labelsep=colon,justification=raggedright,singlelinecheck=false,skip=5pt]{caption}
\DeclareCaptionLabelSeparator{zhcolon}{：}
\captionsetup{labelsep=zhcolon}
\captionsetup[figure]{name=图表,position=top}
\captionsetup[table]{name=图表,position=top}
\makeatletter\let\c@table\c@figure\makeatother
% 图表下方一行资料来源（唯一允许的小字）
\newcommand{\source}[1]{\par\vspace{2pt}{\raggedright\sffamily\footnotesize\color{rgrey}资料来源：#1\par}}

\usepackage{booktabs,tabularx,array,multirow}
\renewcommand{\arraystretch}{1.3}
\newcolumntype{L}{>{\raggedright\arraybackslash}X}
\newcolumntype{C}{>{\centering\arraybackslash}X}

\usepackage{amsmath,amssymb,bm}
\usepackage{enumitem}
\setlist{nosep,leftmargin=2em,itemsep=0.2em}

\usepackage{tikz}
\usetikzlibrary{arrows.meta,positioning,calc,fit,decorations.pathreplacing}
\tikzset{
  node/.style = {draw=rnavy, line width=0.6pt, rounded corners=1.5pt, align=flush center,
                 inner xsep=6pt, inner ysep=4pt, font=\sffamily\small},
  key/.style  = {node, draw=raccent, line width=0.8pt},
  soft/.style = {node, draw=rline},
  flow/.style = {-{Stealth[length=2mm]}, line width=0.6pt, draw=rnavy},
  note/.style = {font=\sffamily\small, text=rgrey, align=flush center},
}

% 首页研究要点：每条以一句加粗结论开头
\newenvironment{keypoints}{%
  \par\noindent{\sffamily\large\bfseries\color{rnavy}研究要点}\par\vspace{0.3em}
  {\color{rnavy}\hrule height 0.6pt}\vspace{0.6em}
  \begin{itemize}[leftmargin=1.2em,itemsep=0.6em,label={\color{rnavy}\rule[0.3ex]{0.45em}{0.45em}}]}
  {\end{itemize}\vspace{0.4em}{\color{rline}\hrule height 0.4pt}}
\newcommand{\kp}[1]{{\sffamily\bfseries\color{rnavy}#1}}

\usepackage[hidelinks]{hyperref}
\usepackage{bookmark}

% 封面
\makeatletter
\newcommand{\subtitle}[1]{\gdef\@subtitle{#1}}
\newcommand{\team}[1]{\gdef\@team{#1}}
\newcommand{\reportdate}[1]{\gdef\@reportdate{#1}}
\providecommand{\@subtitle}{}\providecommand{\@team}{}\providecommand{\@reportdate}{}
\newcommand{\makereporttitle}{%
  \thispagestyle{fancy}
  \vspace*{0.4em}
  {\raggedright\sffamily\color{rnavy}\fontsize{20pt}{27pt}\selectfont\bfseries\@title\par}
  \vspace{0.5em}
  {\raggedright\sffamily\large\color{rnavy}\@subtitle\par}
  \vspace{0.3em}
  {\raggedright\sffamily\small\color{rgrey}\@team\qquad\@reportdate\par}
  \vspace{1.2em}}
\makeatother

%   \title{…} \subtitle{…} \team{…} \reportdate{…}
%   \begin{document} \makereporttitle ...

\setCJKmainfont{Songti SC}[BoldFont={Songti SC Bold}]
\def\yhpath{/Applications/Microsoft PowerPoint.app/Contents/Resources/DFonts/}
\setCJKsansfont{msyh.ttc}[Path=\yhpath,BoldFont=msyhbd.ttc]
\setCJKmonofont{msyh.ttc}[Path=\yhpath]
\setmainfont{Times New Roman}
\setsansfont{msyh.ttc}[Path=\yhpath,BoldFont=msyhbd.ttc]
\setmonofont{Menlo}[Scale=0.88]

\usepackage{newunicodechar}
\newunicodechar{→}{\ensuremath{\rightarrow}}
\newunicodechar{≈}{\ensuremath{\approx}}
\newunicodechar{≤}{\ensuremath{\leq}}
\newunicodechar{≥}{\ensuremath{\geq}}
\newunicodechar{−}{\ensuremath{-}}
\newunicodechar{×}{\ensuremath{\times}}
\newunicodechar{÷}{\ensuremath{\div}}
\newunicodechar{±}{\ensuremath{\pm}}
\newunicodechar{β}{\ensuremath{\beta}}
\newunicodechar{λ}{\ensuremath{\lambda}}
\newunicodechar{σ}{\ensuremath{\sigma}}

\usepackage{xcolor}
\definecolor{rnavy}{HTML}{1B365D}
\definecolor{rgrey}{HTML}{7A828C}
\definecolor{rline}{HTML}{BFC5CC}
\definecolor{raccent}{HTML}{C0392B}

\usepackage[a4paper,top=2.6cm,bottom=2.6cm,left=2.7cm,right=2.7cm,headheight=14pt]{geometry}
\linespread{1.45}
\setlength{\parskip}{0.35em}

% 标题层级：章（\section）、节（\subsection）、小节（\subsubsection，少用）
\ctexset{
  section={format=\sffamily\Large\bfseries\color{rnavy}\raggedright, name={,}, number=\arabic{section},
           aftername=\enspace, beforeskip=1.6em, afterskip=0.8em,
           numberformat={\color{white}\colorbox{rnavy}{\makebox[1.1em]{\arabic{section}}}}},
  subsection={format=\sffamily\large\bfseries\color{rnavy}\raggedright, beforeskip=1.2em, afterskip=0.5em},
  subsubsection={format=\sffamily\normalsize\bfseries\raggedright, beforeskip=0.9em, afterskip=0.3em},
}

% 页眉：左栏目、右日期，下方细色线；页脚：页码
\usepackage{fancyhdr}
\providecommand{\ReportHeadLeft}{策略研究报告}
\providecommand{\ReportHeadRight}{2026 年 11 月}
\pagestyle{fancy}
\fancyhf{}
\fancyhead[L]{\sffamily\small\color{rnavy}\ReportHeadLeft}
\fancyhead[R]{\sffamily\small\color{rnavy}\ReportHeadRight}
\renewcommand{\headrule}{{\color{rnavy}\hrule height 0.8pt}}
\fancyfoot[C]{\sffamily\small\thepage}
\fancypagestyle{plain}{\fancyhf{}\renewcommand{\headrulewidth}{0pt}\fancyfoot[C]{\sffamily\small\thepage}}

% 图表：标题在上，「图 1　标题」
\usepackage{graphicx,float}
\usepackage[font=small,labelfont={sf,bf,color=rnavy},textfont={sf,bf,color=rnavy},labelsep=colon,justification=raggedright,singlelinecheck=false,skip=5pt]{caption}
\DeclareCaptionLabelSeparator{zhcolon}{：}
\captionsetup{labelsep=zhcolon}
\captionsetup[figure]{name=图表,position=top}
\captionsetup[table]{name=图表,position=top}
\makeatletter\let\c@table\c@figure\makeatother
% 图表下方一行资料来源（唯一允许的小字）
\newcommand{\source}[1]{\par\vspace{2pt}{\raggedright\sffamily\footnotesize\color{rgrey}资料来源：#1\par}}

\usepackage{booktabs,tabularx,array,multirow}
\renewcommand{\arraystretch}{1.3}
\newcolumntype{L}{>{\raggedright\arraybackslash}X}
\newcolumntype{C}{>{\centering\arraybackslash}X}

\usepackage{amsmath,amssymb,bm}
\usepackage{enumitem}
\setlist{nosep,leftmargin=2em,itemsep=0.2em}

\usepackage{tikz}
\usetikzlibrary{arrows.meta,positioning,calc,fit,decorations.pathreplacing}
\tikzset{
  node/.style = {draw=rnavy, line width=0.6pt, rounded corners=1.5pt, align=flush center,
                 inner xsep=6pt, inner ysep=4pt, font=\sffamily\small},
  key/.style  = {node, draw=raccent, line width=0.8pt},
  soft/.style = {node, draw=rline},
  flow/.style = {-{Stealth[length=2mm]}, line width=0.6pt, draw=rnavy},
  note/.style = {font=\sffamily\small, text=rgrey, align=flush center},
}

% 首页研究要点：每条以一句加粗结论开头
\newenvironment{keypoints}{%
  \par\noindent{\sffamily\large\bfseries\color{rnavy}研究要点}\par\vspace{0.3em}
  {\color{rnavy}\hrule height 0.6pt}\vspace{0.6em}
  \begin{itemize}[leftmargin=1.2em,itemsep=0.6em,label={\color{rnavy}\rule[0.3ex]{0.45em}{0.45em}}]}
  {\end{itemize}\vspace{0.4em}{\color{rline}\hrule height 0.4pt}}
\newcommand{\kp}[1]{{\sffamily\bfseries\color{rnavy}#1}}

\usepackage[hidelinks]{hyperref}
\usepackage{bookmark}

% 封面
\makeatletter
\newcommand{\subtitle}[1]{\gdef\@subtitle{#1}}
\newcommand{\team}[1]{\gdef\@team{#1}}
\newcommand{\reportdate}[1]{\gdef\@reportdate{#1}}
\providecommand{\@subtitle}{}\providecommand{\@team}{}\providecommand{\@reportdate}{}
\newcommand{\makereporttitle}{%
  \thispagestyle{fancy}
  \vspace*{0.4em}
  {\raggedright\sffamily\color{rnavy}\fontsize{20pt}{27pt}\selectfont\bfseries\@title\par}
  \vspace{0.5em}
  {\raggedright\sffamily\large\color{rnavy}\@subtitle\par}
  \vspace{0.3em}
  {\raggedright\sffamily\small\color{rgrey}\@team\qquad\@reportdate\par}
  \vspace{1.2em}}
\makeatother
